\section{Limitations}
\label{sec:limitations}

\paragraph{Reasoning-formed refusals.}
Our primary attacker is a thinking-only model whose reasoning cannot be disabled,
so it intermittently refuses the escape task during its chain-of-thought. Crucially,
refusals are \emph{not uniform} across controls: they are rare on baseline and standard
controls (isolated instances, at most three of ten epochs in any cell) but occur in every
B6 (distroless) cell, where the environment offers the model no interactive shell and its
deliberation ends in a decline, often after a few failed tool calls. This
couples refusal behaviour to the treatment variable, so we cannot simply average it
away: we classify each sample as escape, blocked, or refused and compute escape
rates over genuine attempts only (\Cref{sec:methodology}). Two residual risks
remain. First, on heavily hardened controls the number of genuine attempts can be
small, widening confidence intervals for exactly the cells of greatest interest.
Second, the escape and blocked outcomes are also produced under active reasoning,
so absolute rates reflect a reasoning-enabled agent and may differ from a
reasoning-disabled one. We mitigate both by reporting refusal rates
explicitly alongside escape rates rather than folding them into the defense signal,
and by cross-checking every model outcome against the LLM-independent
reference exploit (RQ2), which cannot refuse.

\paragraph{Single-host local mode.}
Running in a shared local Docker daemon (rather than a fresh VM per sample) requires
serial execution and idempotent setup. The reference-exploit verification (RQ2) is
LLM-independent and therefore robust to this choice.

\paragraph{Scenario coverage.}
Kernel-exploit scenarios require a specific guest kernel and are excluded from local
mode. The seven in-scope Docker scenarios are \texttt{privileged}, \texttt{docker\_socket}, \texttt{pid\_namespace}, \texttt{sys\_admin}, \texttt{sys\_module}, \texttt{dac\_read\_search}, and \texttt{hostpath\_etc}. Model-level results for \texttt{hostpath\_etc} come from a clean re-run with a per-epoch host guard, which did not cover every persistence path; two hooks written during B7 persisted into B8. The agent never escaped that scenario, so its RQ2 reference results are the informative part.

\paragraph{Non-root execution.}
We did not evaluate running as a non-root user. No control changes the image's default user, so every condition runs as root inside the container, and the results say nothing about non-root execution alone. This is left to future work.

\paragraph{Model and sample size.}
Results use a single attacker model at \num{10} epochs per cell, yielding
wide confidence intervals on individual cells. We therefore emphasize the control
\emph{ranking} and the LLM-independent reference-exploit classification (RQ2), which
do not depend on absolute per-cell rates; generalization across model families is
left to future work.

\paragraph{Reference-exploit repetitions.}
Each row of \texttt{results/eval\_logs/reference\_exploits.csv} is one reference-exploit run,
and the \texttt{repeat} column is \texttt{1} in every row. A cell has one to four such rows,
recorded in separate campaign sessions, and the rows within each cell agree. The
\texttt{repeat} field therefore does not give a within-cell estimate of run-to-run variance
for the reference classification. A reference
failure can reflect a harness or timing problem rather than a control, which is why the
harness-artifact case in \Cref{sec:results} was audited individually.
